\renewcommand{\thetable}{3.A.9}
\begin{tablalafrox}{Registro de decisiones de alcance y diseño. Fuente: elaboración propia; Bases y Subdocumentos 1 y 2 según se indica.}%
  {tab:3-A-9}%
  {F{0.060} F{0.118} Y F{0.300}}%
  {\cab{ID} & \cab{Materia} & \cab{Decisión} & \cab{Fundamento}}
  D-01 & Reparto entre etapas & Etapa 1: trazabilidad, frío, preventa, entrega, efectivo, rutas y OTIF. Etapa 2: canal moderno, portal y costo de servir. & La segunda etapa depende de hechos de entrega registrados; las rutas se adelantan por la jubilación del planificador (LafroX, 2026a, Anexo 2.1, F-06 y F-14; Anexo 2.2, S-16). \\
  D-02 & Medición de OTIF & M10 calcula diariamente el OTIF desde la Etapa 1 conforme a RNG-09, utilizando una única definición de cumplimiento, contando cada pedido una vez y conservando el historial de reentregas y cambios de promesa. & RNG-09; (LafroX, 2026a, Anexo 2.2, S-01); RF-11.03 y R18-04. Evita definiciones distintas entre áreas y mantiene trazabilidad del cálculo. \\
  D-03 & Promesa de entrega & La promesa de entrega es de 24 horas para clientes urbanos con pedidos ingresados antes de las 14:00 y de 48 horas para los urbanos ingresados a las 14:00 o después y para clientes rurales, periféricos o abastecidos mediante cross-docking. Comercial y Operaciones son responsables de esta regla de corte y promesa, que se muestra en preventa. & (LafroX, 2026c, sección «Arbitraje de tensiones operacionales y comerciales», arbitraje 1; Anexo 3.C, S-21). La promesa segmentada evita comprometer un plazo universal y permite informar al cliente las condiciones de entrega. \\
  D-04 & Aceptación & Los resultados de negocio comprometidos se verifican con evidencia objetiva y responsables identificados, en las etapas y períodos establecidos para cada criterio de aceptación. El cierre de cada marcha blanca exige cumplir las condiciones de las Bases Administrativas y contar con el acta firmada por la Contraparte Técnica. Los resultados sujetos a medición posterior se comprueban al completar el período correspondiente. & \textit{Bases Administrativas} (2026, arts. 17.3 y 18.2, p. 13). \\
  D-05 & Reposición & La sugerencia de reposición considera venta real, estacionalidad, promociones comprometidas y plazo de cada proveedor. El Jefe de Abastecimiento la revisa y aprueba antes de su traspaso al ERP, que conserva la emisión de las órdenes de compra. Toda modificación de la sugerencia registra responsable y motivo. & (LafroX, 2026a, Tabla 2.A.8); RF-10.01 a RF-10.03. Sustituye la planilla de reposición y mantiene la aprobación humana antes de generar la compra. \\
  D-06 & Implantación & La implantación se realiza gradualmente por proceso, sitio o zona, con usuarios capacitados y conciliación diaria durante la convivencia con la operación vigente. El Gerente de Operaciones valida el avance de cada ola según sus criterios documentados. & \textit{Bases Administrativas} (2026, art. 17.3, p. 13); \textit{Bases Técnicas Transversales} (2026, cap. 20, RT-20.02--RT-20.04, p. 35). El avance exige evidencia de preparación operacional y aceptación del alcance correspondiente. \\
  D-07 & Reversión & El responsable operativo del CLIENTE autoriza la reversión y el Jefe de Proyecto de LafroX coordina su ejecución. Los registros capturados se conservan para su conciliación posterior. & \textit{Bases Técnicas Transversales} (2026, cap. 20, RT-20.02, p. 35); (LafroX, 2026b, sección «Estructura para Proyecto»); (LafroX, 2026a, Anexo 2.2, S-14). Preserva la evidencia y evita pérdida o duplicación. \\
  D-08 & Modo desconectado & Terreno registra la jornada íntegra durante 14 h sin cobertura; el CD recibe, prepara y despacha durante 24 h sin enlace. Captura local, identificador único y sincronización idempotente: máximo 10 min para reparto y 2 h para CD tras reconectar. & \textit{Bases Técnicas Transversales} (2026, cap. 3, RT-03.12, p. 9); (LafroX, 2026a, Anexo 2.2, S-08): stock local indicativo, reserva al confirmar en el servidor y conflictos resueltos al sincronizar con notificación del quiebre. \\
  D-09 & Analítica avanzada & Ruteo determinístico y explicable; la solución no incorpora componentes de inteligencia artificial. & El planificador debe revisar y corregir la ruta (R18-07; LafroX, 2026a, Anexo 2.1, F-06; Anexo 2.2, S-16). La incorporación de IA no es obligatoria (\textit{Bases Técnicas Transversales}, 2026, cap. 18, párrafo introductorio, p. 32). \\
  D-10 & Mesa de servicio & Los incidentes críticos y altos se atienden 24$\times$7$\times$365, con tiempos máximos de respuesta de 15 minutos y 1 hora, y de resolución de 4 y 8 horas, respectivamente. La atención general a usuarios opera de 04:00 a 22:00 de lunes a sábado y 24$\times$7 en septiembre y diciembre. Cada incidente se registra en un ticket único, se clasifica por severidad y se escala entre los niveles L1, L2 y L3 según corresponda, con seguimiento hasta su cierre. & \textit{Bases Administrativas} (2026, art. 78.2, p. 40); \textit{Bases Administrativas} (2026, art. 79, pp. 40--41); \textit{Bases Técnicas Transversales} (2026, cap. 21, RT-21.15, p. 37); (LafroX, 2026b, sección «Centro de Operaciones y Mesa de Servicios»). \\
  D-11 & Ambientes & Cinco ambientes: DEV, QA, PREPROD, PROD y DR, habilitados y operativos para H3. Capacitación usa un espacio aislado con datos ficticios; no constituye un sexto ambiente. & (LafroX, 2026b, sección «Presentación de la empresa»); \textit{Bases Técnicas Transversales} (2026, cap. 4, RT-04.01, p. 10). \\
  D-12 & Cobertura de pruebas & Umbral bloqueante de 80 \% de cobertura de pruebas unitarias. & Política del Subdocumento 1 (LafroX, 2026b, sección 1.3.1 «Modelo de Gobierno de Calidad», dentro de 1.3 «Gobierno interno Calidad, Seguridad y Conocimiento»), que supera el 70 \% de \textit{Bases Técnicas Transversales} (2026, cap. 4, RT-04.11, p. 11). \\
\end{tablalafrox}
